\todo[inline,color=gray!10]{meta: Dependencies: Sections 04 (what we do), 05 (what we demonstrate).}

\todo[inline,color=gray!10]{meta: Note: Double-blind. CSCS 2024 short paper referenced in third person.}

\todo[inline,color=gray!10]{meta: Source: Gap analysis synthesis (2026-03-25), archive/gap analyis/GAP\_ANALYSIS\_SYNTHESIS.md}

\section{Related Work}\label{related-work}

\label{sec:related-work}

\subsection{Credential Ecosystem Design and Formalization}\label{credential-ecosystem-design-and-formalization}

\label{sec:rw-credential}

Verifiable credential schema design is currently guided by specifications that define credential structure at individual abstraction layers: the W3C Verifiable Credentials Data Model 2.0 \citep{manu_verifiable_2025} and EU Architecture Reference Framework \citep{noauthor_eu-digital-identity-walleteudi-doc-architecture-and-reference-framework_2026} prescribe issuance and presentation flows, while AnonCreds \citep{curran2022anoncreds} and ISO mDL \citep{_mobile_2021} define format-specific encoding and proof mechanisms. \cutcandidate[supplementary survey reference — can be removed for budget]{Mazzocca et al. \citep{mazzocca_survey_2025} provide a comprehensive survey of this landscape.} None of these specifications provides a formal mechanism for checking constraints that span multiple layers. Where formal methods have been applied, they target individual layers: Braun and Käfer \citep{curry_rdf-based_2025} define RDF-based semantics for selective disclosure and zero-knowledge proofs on verifiable credentials, Yamamoto et al. \citep{yamamoto_formalising_2022} formalize selective disclosure specifically for linked-data credentials, and Braun et al. \citep{braun_ssi_2024} verify SSI protocol security properties using ProVerif. Each formalization targets protocol-level security or single-format semantics --- none operates across the boundary between domain-level claim semantics and format-specific representation capabilities.

Complementary conceptual models address individual concerns in credential system design. Tith and Colin \citep{tith_trust_2025} propose a trust policy meta-model that captures how identity systems establish and evaluate trust across organizational boundaries. Turkanović et al. \citep{turkanovic_model_2025} formalize delegated authority enforcement, modeling how credential issuance rights propagate through delegation chains. The Trust over IP stack \citep{davie_trust_2019} organizes the credential ecosystem into informal governance layers, and Naghmouchi and Laurent \citep{naghmouchi_systematic_2025} systematize privacy-by-design principles for SSI. \cutcandidate[supplementary conceptual models — can be removed for budget]{Garcia-Rodriguez et al. \citep{garcia-rodriguez_towards_2021} propose a standardized model for privacy-preserving credentials, and Schardong and Custódio \citep{maass_role-artifact-function_2025} develop a framework for understanding digital identity models.} Each of these works addresses a single concern --- trust policy, delegation, governance layering, or privacy --- but none formalizes cross-layer constraints spanning domain semantics, credential structure, and format-specific representation simultaneously.

\subsection{Model-Driven Engineering for Security and SSI}\label{model-driven-engineering-for-security-and-ssi}

\label{sec:rw-mde}

Model-driven security engineering has established metamodel-based approaches to security properties of software architectures. UMLsec \citep{noauthor_secure_2005} annotates UML models with confidentiality and authentication constraints, enabling formal verification of security properties during design. SecureUML \citep{basin_model_2006} integrates role-based access control specifications into class models, generating enforcement infrastructure from the model. Neither targets credential schema design --- both operate on software architecture elements rather than the domain-specific structure of verifiable credentials.

Within SSI specifically, four model-driven approaches address different facets of the design problem. Cippitelli et al. \citep{cippitelli_chorssi_2024} model SSI interactions as BPMN choreographies and generate executable code from the choreography models --- their concern is the protocol execution flow between participants, not the structure of credentials exchanged. Ding and Sato \citep{ding_model-driven_2023} apply model-driven security analysis to SSI architectural patterns, systematically identifying threats in system configurations rather than checking credential design consistency. Pattiyanon et al. \citep{pattiyanonMethodDetectingCommon2022} define domain-specific modeling languages and knowledge graphs that surface common SSI implementation weaknesses, while Barclay et al. \citep{barclay_towards_2020} model SSI governance requirements using iStar goal models, operating at the requirements level without formalizing credential structure. In adjacent work, King et al. \citep{king_automated_2017} automate multi-level governance compliance checking using institutional norms --- a structurally related problem, though their formal objects (normative rules governing agent actions) differ from the structural constraints over model elements addressed here.

Model-driven engineering has thus been applied to SSI for choreography, security analysis, weakness detection, and governance requirements. None of these works defines a multi-level metamodel or formalizes cross-layer constraints connecting domain-level semantics through credential structure to format-specific representation.

\subsection{Multi-Level Modeling and Graph-Based Design Space Exploration}\label{multi-level-modeling-and-graph-based-design-space-exploration}

\label{sec:rw-multilevel}

Multi-level metamodeling, as established by Atkinson and Kühne \citep{goos_essence_2001, atkinson_reducing_2008}, eliminates accidental complexity in deep classification scenarios by allowing model elements to span more than two metalevels; de Lara and Guerra \citep{hutchison_deep_2010} implement these principles with deep instantiation in MetaDepth. Diskin et al. \citep{dingel_specifying_2011} formalize consistency checking across heterogeneous model views --- a problem this work extends by adding independently governed constraint sources as a consistency dimension.

The present metamodel relies on the Refinery partial graph modeling framework \citep{marussy_refinery_2024} as its solver infrastructure. Refinery generates diverse model instances that satisfy structural and relational constraints expressed as graph predicates, and its three-valued partial model semantics enable reasoning over designs that are still incomplete --- a property essential for credential ecosystems where not all schema elements are known at design time. A prior short paper \citep{farkas_prolog-based_2024} applied Refinery to credential schema validation with a single-layer prototype; the present work extends this to a three-layer metamodel with formalized cross-layer constraints.

Unlike standard multi-level modeling applications where layers represent successive instantiation and constraints take the form of potency annotations, the three layers in the present metamodel --- domain concepts, credential structure, and format-specific representation --- represent independently governed concern spaces. They are connected by coverage and capability constraints formalized as graph predicates over partial models (\autoref{sec:cross-layer}), not by instantiation relationships. The contribution is therefore the formalization of governance constraints across independently governed layers, not the layering technique itself.

No prior work combines multi-level metamodeling with formalized cross-layer constraints for verifiable credential ecosystem design under multi-source governance.
